% GitHub Repo and Documentation: https://github.com/celiobjunior/resume-template
% Copyright © 2025 Celio B Junior. All rights reserved.
% 
% Licensed under the Apache License, Version 2.0 (the "License");
% you may not use this file except in compliance with the License.
% You may obtain a copy of the License at
%
%     http://www.apache.org/licenses/LICENSE-2.0
%
% This template follows best practices from README.md

% Início do documento LaTeX: define tipo e formato do documento
% a4paper = tamanho A4, 10pt = tamanho base da fonte
\documentclass[a4paper,10pt]{article}

% --- PACOTES: BASE / IDIOMA ---
\usepackage[utf8]{inputenc}
\usepackage[T1]{fontenc}
\usepackage[portuguese]{babel}

% --- PACOTES: LAYOUT / TIPOGRAFIA ---
\usepackage{geometry}
\usepackage{parskip}
\usepackage{microtype}
\usepackage{enumitem}
\usepackage{titlesec}
\usepackage{array}
\usepackage{tabularx}

% --- PACOTES: LINKS / BOOKMARKS ---
\usepackage{hyperref}
\usepackage{bookmark}
\usepackage{xurl}

% --- CONFIGURAÇÃO: MARGENS / ESPAÇAMENTO ---
\geometry{top=1.5cm, bottom=1.5cm, left=1.5cm, right=1.5cm} % Margens da página
\setcounter{secnumdepth}{0}                                 % Remove numeração de seção
\setlist[itemize]{
    leftmargin=0.75em, % Recuo da lista (mais perto da margem da página)
    itemsep=0.2em,     % Espaço entre itens
    topsep=0.25em,     % Espaço antes/depois da lista
    parsep=0em,        % Espaço entre parágrafos no mesmo item
    partopsep=0em      % Espaço extra ao iniciar a lista
}

% Remove números de página e cabeçalhos para visual limpo do currículo
\pagestyle{empty}

% --- CONFIGURAÇÃO: METADADOS / LINKS ---
\hypersetup{
    pdftitle={CV João Pedro Calsavara},   % Título do PDF
    pdfauthor={João Pedro Calsavara},     % Autor do PDF
    colorlinks=true,          % Usa cor nos links (sem caixa)
    linkcolor=black,          % Cor de links internos
    urlcolor=black,           % Cor de URLs
    citecolor=black,          % Cor de citações
    bookmarksdepth=2          % Barra lateral: seção e subseção
}

% --- CONFIGURAÇÃO: TÍTULOS ---
\titleformat{\section}
{\Large\bfseries}
{}
{0em}
{}
[\titlerule\vspace{0.5ex}]
\titleformat{\subsection}
{\normalsize\bfseries}
{}
{0em}
{}
\titlespacing*{\subsection}{0pt}{0.35em}{0.15em}

% Macro para entradas do currículo com bookmark no PDF
\newcounter{cventry}
\newcommand{\cventry}[4]{%
    \refstepcounter{cventry}%
    \phantomsection%
    \pdfbookmark[2]{#1}{cventry-\thecventry}%
    \noindent\begin{tabularx}{\textwidth}{@{}>{\raggedright\arraybackslash}X >{\raggedleft\arraybackslash}X@{}}
    \textbf{#1} & #2 \\
    \textit{#3} & \textit{#4} \\
    \end{tabularx}
}

% --- INÍCIO DO DOCUMENTO ---
\begin{document}

% --- CABEÇALHO ---
% Substitua por suas informações pessoais
\begin{center}
    {\LARGE \textbf{João Pedro Calsavara}} 
    \\ [0.1cm]
    Limeira, SP
    {\textbullet}
    \href{mailto:jpcalsavara@email.com}{jpcalsavara@email.com} 
    {\textbullet}
    \href{https://linkedin.com/in/joaopedrocalsavara}{linkedin.com/in/joaopedrocalsavara} 
    {\textbullet}
    \href{https://github.com/JPCalsavara}{github.com/JPCalsavara}
\end{center}

% --- SEÇÕES ---

\section{Educação}
    % Formação mais recente primeiro
    \cventry{Universidade Estadual de Campinas (UNICAMP)}{Limeira, SP}{Tecnólogo em Análise e Desenvolvimento de Sistemas}{Fev 2023 - Dez 2026}

\section{Experiência}
    \cventry{Mottu}{São Paulo, SP}{Desenvolvedor Backend (Estagiário)}{Set 2025 - Atual}
        \begin{itemize}
            \item Assegurou a consistência de dados operacionais e retenção de repasses financeiros ao desenvolver microsserviços escaláveis em .NET 8 utilizando mensageria (Pub/Sub), Docker e Kubernetes.
            \item Eliminou 2 horas de trabalho manual diário na operação logística arquitetando um CronJob em Kubernetes integrado a LLM (Engenharia de Prompt) para triagem de notificações.
            \item Evitou perdas semanais de penalidades e reduziu o tempo de resolução de incidentes ao criar painéis analíticos e alertas em tempo real no Datadog.
        \end{itemize}

    \cventry{Atria Jr.}{Limeira, SP}{Desenvolvedor Backend / Assessor Comercial}{Mai 2024 - Set 2025}
        \begin{itemize}
            \item Desenvolveu aplicação web backend em TypeScript seguindo Clean Architecture e Domain-Driven Design (DDD), utilizando PostgreSQL, Prisma, Docker e infraestrutura AWS.
            \item Reestruturou estrategicamente o blog com foco em SEO, gerando mais de 1.000 usuários orgânicos mensais para a Empresa Júnior.
            \item Conduziu reuniões diagnósticas atuando como ponte técnica entre cliente e time de desenvolvimento para estruturar propostas comerciais.
        \end{itemize}

\section{Atividades de Liderança}
    \cventry{Semeia Code}{Limeira, SP}{Co-fundador e Coordenador Executivo}{Jul 2024 - Fev 2026}
        \begin{itemize}
            \item Co-fundou projeto de extensão para ampliar o acesso à educação em programação para alunos de escolas públicas no ensino médio.
            \item Coordenou equipes multidisciplinares compostas por alunos da Unicamp, impactando diretamente e continuamente dezenas de estudantes a cada semestre.
        \end{itemize}

\section{Habilidades}
    \begin{itemize}
        \item \textbf{Técnicas:} C#, .NET 8, TypeScript, Python, Next.js, PostgreSQL, SQL Server, Docker, Kubernetes, AWS, RabbitMQ, Datadog, xUnit.
        \item \textbf{Idiomas:} Português (Nativo), Inglês (Avançado).
    \end{itemize}

\end{document}
